\chapter{Infraestructura de Cámaras}
\label{ch:cameras}
% ==============================================================

\section{Sistema Obscape}

Las seis cámaras están gestionadas por la plataforma Obscape y se acceden via API REST
autenticada. La tabla~\ref{tab:cameras} lista los parámetros de registro de cada estación.

\begin{table}[H]
\centering
\caption{Estaciones Obscape en Guardamar del Segura}
\label{tab:cameras}
{\small\resizebox{\textwidth}{!}{%
\begin{tabular}{C{1.2cm} C{1.5cm} C{2.3cm} C{2.8cm} C{2.8cm} C{2cm}}
\toprule
\textbf{Cámara} & \textbf{Station} & \textbf{Referencia} & \textbf{Latitud (°N)} & \textbf{Longitud (°E)} & \textbf{Última act.} \\
\midrule
CAM~1 & 8213 & PTM61471 & 38.110327 & $-$0.643027 & 20-05-2026 \\
CAM~2 & 8214 & PTM61474 & 38.096291 & $-$0.645725 & 20-05-2026 \\
CAM~3 & 8212 & PTM61473 & 38.087348 & $-$0.646994 & 20-05-2026 \\
CAM~4 & 8211 & PTM61475 & 38.087385 & $-$0.647078 & 20-05-2026 \\
CAM~5 & 8209 & PTM61472 & 38.076774 & $-$0.648117 & 20-05-2026 \\
CAM~6 & 8210 & PTM61470 & 38.076796 & $-$0.648114 & 20-05-2026 \\
\bottomrule
\end{tabular}
}}
\end{table}

\begin{warnbox}[Nota geográfica]
CAM~3 y CAM~4 son casi coincidentes ($\Delta \approx 8$~m) y cubren el mismo tramo desde
ángulos ligeramente distintos. Lo mismo ocurre con CAM~5 y CAM~6. Esto permite comparación
cruzada de la línea detectada para validación interna.
\end{warnbox}

\section{Distribución geográfica}

Las cámaras se distribuyen de norte a sur a lo largo de $\approx 3{,}5$~km de frente litoral.
La cobertura solapada garantiza continuidad en la extracción de línea de costa sin lagunas.

\section{Descarga automática de imágenes}

El módulo \texttt{acces\_api/scheduled\_download.py} implementa una \textbf{ventana deslizante de
14 días} que garantiza recuperación automática ante fallos puntuales de descarga:

\begin{codebox}[Descarga estándar]
\begin{minted}{bash}
# Últimas 2 semanas, hora solar 12:00 (por defecto)
python3 acces_api/scheduled_download.py

# Personalizar ventana y hora
python3 acces_api/scheduled_download.py --days 30 --hour 15

# Descarga masiva histórica (todas las horas)
python3 acces_api/scheduled_download.py --days 90 --hour all
\end{minted}
\end{codebox}

El sistema evita duplicados comprobando la existencia del fichero antes de descargar y utiliza
un nombre determinista: \texttt{UNIX\_YYYYMMDD\_HHMMSS\_REF.jpg}.

\subsection{Automatización con Cron}

Para descarga diaria autónoma:

\begin{codebox}[/etc/cron.d/cvlit-download]
\begin{minted}{text}
# Descarga diaria a las 13:00h (12h solar + margen de publicación API)
0 13 * * * cd /home/nickhernd/Desktop/cv-lit && \
  python3 acces_api/scheduled_download.py >> data/logs/cron.log 2>&1
\end{minted}
\end{codebox}

% ==============================================================
